\subsection{Measuring and decomposing the spread}

\begin{qfdefinition}[Spread measures: quoted, inside, relative, effective and realised spreads]\label{def:spread-measures}
\qffield{Statement} With best bid $b_t$, best ask $a_t$ and midpoint $m_t=\tfrac12(a_t+b_t)$:
\begin{itemize}
\item the \emph{quoted spread} is $a_t-b_t$; the \emph{inside spread} is the highest bid's distance from the lowest ask across all dealers, and is no wider than the narrowest dealer's spread \citep[chap.~``Dealers'', p.~280]{alpha-harris};
\item the \emph{relative spread} is the spread as a fraction of price, the only basis on which spreads of different instruments compare \citep[chap.~``Bid/Ask Spreads'', p.~312]{alpha-harris};
\item for a trade at $p_t$ with sign $Q_t=+1$ if buyer-initiated and $-1$ if seller-initiated, the \emph{effective spread} is twice the signed difference from the midpoint at the time of the trade, and the \emph{realised spread} twice the signed difference from the midpoint a set interval $\Delta$ later, commonly 5, 10, 15 or 60 minutes \citep[chap.~``Liquidity and Transaction Cost Measurement'', pp.~424--425]{alpha-harris}:
\end{itemize}
\[ ES_t=2Q_t(p_t-m_t), \qquad RS_t=2Q_t(p_t-m_{t+\Delta}), \qquad ES_t=RS_t+2Q_t(m_{t+\Delta}-m_t). \]
The last term, the signed midpoint revision after the trade, is sometimes used as an estimate of the trade's price impact \citep[sec.~14.2.2, eq.~(14.2), p.~148]{s6-hasbrouck}.

\qffield{Notation} $Q_t$ is the trade sign. Harris's effective and realised spreads are doubled, full-spread amounts; Hasbrouck's \emph{effective cost} $p_t-m_t$ and \emph{realised cost} $p_t-m_{t+5}$ are per-side amounts written for a buy order, with the midpoint five minutes after the order's receipt, as in the SEC's Rule 605 \citep[same sec., p.~148]{s6-hasbrouck}. Harris calls the undoubled signed difference the \emph{liquidity premium} \citep[chap.~``Liquidity and Transaction Cost Measurement'', p.~425]{alpha-harris}. Alternative symbols are $\epsilon_t$ for $Q_t$, $s_t$ for the quoted spread and $PI_t$ for the impact term; the relative spread is measured against $m_t$.

\qffield{Conditions} As of the 2002 and 2007 printings for the reporting rules (Rule 605 named by Hasbrouck); definitions timeless. The effective spread equals the quoted spread when every purchase is at the ask and every sale at the bid; it is smaller with price improvement and larger when large orders fill outside the quotes. Realised and effective spreads are equal if the midpoint does not move over the interval; usually it moves with the trade, so realised spreads are smaller \citep[same chap., p.~425]{alpha-harris}.

When the analyst did not take part in the trades, signs are inferred, usually by the Lee--Ready rule, which overstates costs for traders who do not only take liquidity and cannot see orders filled in several trades \citep[same chap., p.~423]{alpha-harris}. Delayed trade reports bias effective costs down, and hidden orders and market-maker improvement distort the inferred sign \citep[sec.~14.3.3, pp.~150--151]{s6-hasbrouck}.

\qffield{Practice} Quoted minus effective spread measures the price improvement dealers give; effective minus realised spread measures dealers' losses to well-informed traders \citep[chap.~``Liquidity and Transaction Cost Measurement'', pp.~425--426]{alpha-harris}. The realised cost is the revenue of a dealer who sold to the customer at $p_t$ and covered at the later midpoint \citep[sec.~14.2.2, p.~148]{s6-hasbrouck}. The effective spread is the least noisy benchmark estimator, because its benchmark is contemporaneous; the realised spread grows noisier the longer the interval \citep[chap.~``Liquidity and Transaction Cost Measurement'', p.~427]{alpha-harris}. The quoted spread measures the cost of immediacy for small orders only; depth and resiliency are other aspects of liquidity \citep[chap.~``Bid/Ask Spreads'', p.~312]{alpha-harris}.

\qffield{Pitfalls} Harris also uses \emph{realised spread} for the difference between the prices at which a dealer actually sells and buys on a round trip, which can be negative when quotes move against the inventory \citep[chap.~``Dealers'', p.~280]{alpha-harris}; that is a dealer's outcome, not the midpoint-based estimator above. Mixing doubled and per-side conventions changes figures by a factor of two. Interpretation: effective minus realised spread is the permanent component only if the interval is long enough for prices to adjust and short enough to exclude unrelated news, and quote moves for inventory that have not yet reversed fall into it too; see \hyperref[res:spread-component-identification]{component identification}.

\qffield{See also} \hyperref[res:adverse-selection-spread-component]{spread components} (what the measures estimate); \hyperref[res:roll-spread-estimator]{Roll's estimator} (alternative: a spread from trade prices alone); \hyperref[def:outside-spread]{outside spread} (contrast: an undisplayed willingness to trade).

\end{qfdefinition}

\begin{qfresult}[Roll's serial-covariance spread estimator]\label{res:roll-spread-estimator}
\qffield{Statement} Let trade prices be $P_t=V_t+\tfrac12SQ_t$, where the fundamental value $V_t$ is a random walk with independent innovations $\varepsilon_t$ of mean zero and variance $\sigma^2$, $S$ is the spread, and the trade sign $Q_t=\pm1$ is a buy with probability one-half, independently of earlier signs and of $\varepsilon_t$ and $\varepsilon_{t+1}$. Then \citep[chap.~``Volatility'', box, p.~415]{alpha-harris}
\[ \operatorname{Var}(\Delta P_t)=\sigma^2+\tfrac12S^2, \qquad \operatorname{Cov}(\Delta P_t,\Delta P_{t-1})=-\tfrac14S^2, \qquad S=2\sqrt{-\operatorname{Cov}(\Delta P_t,\Delta P_{t-1})}, \]
and all autocovariances of order two or more are zero \citep[sec.~3.4, p.~30]{s6-hasbrouck}. Roll's estimator replaces the covariance by its sample value \citep[chap.~``Volatility'', p.~415]{alpha-harris}.

\qffield{Notation} $S$ is the full spread, in price units; changes are between successive trades. Hasbrouck writes $p_t=m_t+cq_t$ with half-spread $c=S/2$, efficient-price innovation $u_t$ of variance $\sigma_u^2$, and $\gamma_0,\gamma_1$ for the variance and first autocovariance of price changes, so $\gamma_0=2c^2+\sigma_u^2$, $\gamma_1=-c^2$, $c=\sqrt{-\gamma_1}$ and $\sigma_u^2=\gamma_0+2\gamma_1$ \citep[sec.~3.4, eqs.~(3.3)--(3.5), pp.~29--30]{s6-hasbrouck}. Harris calls the estimate an effective spread \citep[chap.~``Liquidity and Transaction Cost Measurement'', p.~434]{alpha-harris}.

\qffield{Conditions} The spread is constant and noninformational: competitive dealers quote $m_t\pm c$ to cover a cost $c$ per trade \citep[sec.~3.4, p.~29]{s6-hasbrouck}. Idea of the proof (full derivation: \qfderivation{drv:roll-spread-estimator}): $\Delta P_t=\varepsilon_t+\tfrac12S(Q_t-Q_{t-1})$; every cross-product has zero expectation except $\mathbb{E}Q_{t-1}^2=1$, which gives $-\tfrac14S^2$ for the first autocovariance \citep[chap.~``Volatility'', p.~415]{alpha-harris}. The variance splits into a fundamental part $\sigma^2$ and a transitory part from bid--ask bounce \citep[same chap., pp.~415--416]{alpha-harris}.

\qffield{Practice} The estimator is used where quotes are unavailable \citep[sec.~3.4, p.~30]{s6-hasbrouck}; Harris advises using it only then \citep[chap.~``Liquidity and Transaction Cost Measurement'', p.~434]{alpha-harris}. For one NYSE stock (PCO, October 2003), $\hat\gamma_1=-0.0000294$ gives $c=\$0.017$ and a spread of $\$0.034$, against a time-weighted average quoted spread of $\$0.032$ \citep[sec.~3.4, p.~30]{s6-hasbrouck}.

\qffield{Pitfalls} In that sample the assumptions fail: the spread ranged from \$0.01 to \$0.49, successive trade signs had correlation about $0.34$, and midpoint changes were positively correlated with the latest trade sign \citep[sec.~3.4, p.~30]{s6-hasbrouck}. The biases go both ways. Positively autocorrelated signs, $\operatorname{Corr}(q_t,q_{t-1})=\rho>0$, give $\gamma_1=-c^2(1-2\rho)$ and bias $\hat c$ down; signs positively correlated with the efficient-price innovation, as with informed trading, give $\gamma_1=-c(c+\rho\sigma_u)$ and bias $\hat c$ up \citep[sec.~4.4, Exercises~4.2--4.3, pp.~38--39]{s6-hasbrouck}. A positive sample autocovariance is not uncommon, and estimation error alone can produce one when the model is correct; the square root is then undefined \citep[sec.~4.5, p.~39]{s6-hasbrouck}. Harris says violations typically make it underestimate the spread \citep[chap.~``Liquidity and Transaction Cost Measurement'', p.~434]{alpha-harris}, and that reversals lasting beyond one trade make Roll's variance components understate transitory volatility \citep[chap.~``Volatility'', p.~416]{alpha-harris}.

One classification lists informed trading among causes of a less negative covariance; under Hasbrouck's Exercise~4.3 it makes the covariance more negative. Price discreteness and hidden liquidity, also listed there, are not discussed in the source.

\qffield{See also} \hyperref[res:spread-component-identification]{component identification} (generalisation: the generalized Roll model, with an information component); \hyperref[def:volatility-components]{fundamental and transitory volatility} (generalisation: random-walk decompositions with a general pricing error); \hyperref[def:spread-measures]{spread measures} (when quotes are available); \hyperref[res:grossman-miller-liquidity-measures]{Grossman--Miller liquidity measures} (contrast: negative autocovariance from inventory risk).


\end{qfresult}

\begin{qfresult}[Separating transitory and permanent spread components: bounce, revisions and identification]\label{res:spread-component-identification}
\qffield{Statement} The transaction-cost component makes prices bounce between bid and ask, so its price changes reverse; the adverse-selection component moves the dealers' value estimates, so its price changes are permanent and unpredictable; both depend on whether a buyer or a seller initiated the trade, and the usual decomposition explains current and future price changes with the signs of trades \citep[chap.~``Bid/Ask Spreads'', pp.~299--300 and~302--303]{alpha-harris}. In Hasbrouck's generalized Roll model the efficient price moves with the trade and with public news, and trades occur at a cost $c$ either side of it \citep[sec.~8.2, eq.~(8.1), pp.~67--68]{s6-hasbrouck}:
\[ m_t=m_{t-1}+\lambda q_t+u_t, \qquad p_t=m_t+cq_t, \qquad \text{ask}-\text{bid}=2(c+\lambda), \]
\[ \gamma_0=c^2+(c+\lambda)^2+\sigma_u^2, \qquad \gamma_1=-c(c+\lambda), \]
with all higher autocovariances zero \citep[sec.~8.3, eq.~(8.3), p.~68]{s6-hasbrouck}. Two moments cannot identify the three parameters $\lambda$, $c$, $\sigma_u^2$: trade prices alone do not separate the components. Only the variance of efficient-price changes, $\lambda^2+\sigma_u^2=\gamma_0+2\gamma_1$, is identified without restrictions \citep[same sec., pp.~68--69]{s6-hasbrouck}.

\qffield{Notation} $q_t=\pm1$ is the trade sign; $\lambda$ is the adverse-selection (information) cost and $c$ the noninformational cost, both per side; $u_t$ is nontrade public information, uncorrelated with $q_t$; $\gamma_0,\gamma_1$ are the variance and first autocovariance of $\Delta p_t$. In Harris's terms the adverse-selection component is $2\lambda$ and the transaction-cost component $2c$ (Interpretation).

\qffield{Conditions} Timing: after the trade at $t-1$ the efficient price is $m_{t-1}$; public news $u_t$ arrives; dealers set bid and ask symmetrically about $m_{t-1}+u_t$; the trade $q_t$ arrives and the efficient price becomes $m_t$ \citep[sec.~8.2, p.~68]{s6-hasbrouck}. Setting $\lambda=0$ gives Roll's model; setting $\sigma_u^2=0$ or $c=0$ identifies the rest, but such restrictions are economically strong \citep[sec.~8.3, p.~69]{s6-hasbrouck}.

\qffield{Practice} With trade signs, analyses indicate that in most markets the adverse-selection component is larger than the transaction-cost component \citep[chap.~``Bid/Ask Spreads'', p.~303]{alpha-harris}. Inventory control predicts the same short-run quote response as adverse selection: after a customer buys, the dealer's quotes rise either way. The two separate in the long run, since information effects are permanent while inventory effects vanish once offsetting trades restore the dealer's position; for NYSE specialists, most quote dynamics are attributable to trades, with inventories adding little explanatory power \citep[sec.~11.4, p.~115]{s6-hasbrouck}. Identification therefore needs more than prints: trade signs and subsequent quote revisions and, to separate inventory, dealer positions.

\qffield{Pitfalls} An estimated permanent component includes any inventory-driven quote change that has not reversed within the horizon studied, and the transitory component includes any mean reversion from uninformed order imbalances (Interpretation). Where the dealer is counterparty to every one-unit trade, $q_t=-\Delta I_t$ for inventory $I_t$ \citep[sec.~11.3.2, p.~112]{s6-hasbrouck}, and a regression on current and lagged trades and inventories is then collinear; the NYSE data avoid this because the specialist does not take part in every trade \citep[sec.~11.4, p.~115]{s6-hasbrouck}.

\qffield{See also} \hyperref[res:roll-spread-estimator]{Roll's estimator} (special case: $\lambda=0$); \hyperref[res:adverse-selection-spread-component]{spread components} (the theory being estimated); \hyperref[def:spread-measures]{spread measures} (a model-free split by effective and realised spreads); \hyperref[mod:inventory-based-quotes]{inventory-based quotes} (the inventory alternative); \hyperref[mod:propagator]{propagator model} (impact that decays over many trades rather than one); \hyperref[meth:order-flow-impact-regression]{order-flow impact regression} (the Glosten--Harris extension with trade size).

\end{qfresult}

\begin{qfresult}[Determinants of bid--ask spreads: primary factors, order-driven equilibrium and observable proxies]\label{res:spread-determinants}
\qffield{Statement} Three primary factors set spreads \citep[chap.~``Bid/Ask Spreads'', sec.~14.6.1, pp.~311--313]{alpha-harris}:
\begin{itemize}
\item \emph{asymmetric information} widens them, most when informed traders know material information;
\item \emph{volatility} widens them: it raises the value of the timing option that standing limit orders and dealer quotes give away when they cannot be adjusted at once, makes diversifiable inventory risk costly to risk-averse dealers, and proxies for asymmetric information;
\item \emph{utilitarian trading interest} narrows them: active markets spread dealers' fixed costs over more trades, make inventory easier to lay off, draw public limit orders, and dilute informed flow.
\end{itemize}
In a continuous order-driven market the equilibrium spread makes the marginal trader indifferent between limit and market orders; it would be zero if traders were identical and equally informed, could cancel instantly and free, and cared neither for time nor risk \citep[same chap., secs.~14.4.1--14.4.2, pp.~304--310]{alpha-harris}. Relaxing those assumptions gives Harris's Table~14-2 \citep[same chap., p.~311]{alpha-harris}: information asymmetry (increases spreads; great importance), time to cancel limit orders (increases; high), volatility (increases; high), limit-order management costs (increases; moderate), the value of traders' time (increases; moderate), the commission difference between limit and market orders (one for one), and risk aversion (increases; low).

\qffield{Conditions} Timeless (theory); institutional examples from 2002; the European evidence from 2008--2016. With costly limit-order management, the equilibrium spread equals the expected cost of managing limit orders plus twice the value of the timing option \citep[same chap., p.~307]{alpha-harris}; with only a commission difference, it equals that difference \citep[same chap., p.~306]{alpha-harris}. If all traders were equally uninformed, uncertainty would narrow spreads, because value estimates are less volatile than values; uncertainty widens spreads through asymmetry \citep[same chap., p.~310]{alpha-harris}. The primary factors interact: asymmetry widens spreads, which deters uninformed traders and lowers volume, which widens them further \citep[same chap., p.~312]{alpha-harris}. Public limit-order traders without dealers' costs can undercut dealers, who survive by exploiting their view of the order flow \citep[same chap., sec.~14.5, pp.~310--311]{alpha-harris}.

\qffield{Practice} Observable \emph{secondary} determinants stand in for the primary ones \citep[same chap., sec.~14.6.2, pp.~314--317]{alpha-harris}. For asymmetric information: disclosure rules, market-condition reports, analyst and press coverage, economy-wide commodity contracts, diversified portfolios and diversified firms, established rather than emerging industries, firm age, enforced insider-trading rules, and expected material news such as earnings. For volatility: financial and operating leverage and uncertain growth for stocks; perishability, storage costs, low inventories and weather for commodities; political, inflation and interest-rate uncertainty for currencies; inflation, rates and credit quality for bonds. For utilitarian interest: trading activity (volume, number of trades), firm size, debt issue size and on-the-run status, and contracts that replicate risks many traders want to hedge.

When asymmetric information is severe or utilitarian interest small, spreads can be so wide that the market fails; small firms then finance through banks and venture capitalists, whose due diligence, restrictions and supervision rights need a private relationship \citep[same chap., sec.~14.6.3, p.~317]{alpha-harris}. Empirically, daily spreads rise with daily volatility, the spread rewarding market makers for the risk they bear, and across small-tick instruments spreads follow $\psi_i=a+b\,\sigma_i/\sqrt{N_i}$, with $N_i$ the daily number of trades \citep[chap.~2, secs.~2.2 and~2.3.1, pp.~144--145 and~150--151]{alpha-laruelle}; see \hyperref[meth:spread-volatility-regression]{spread--volatility regression}.

When the spread is stuck at one tick, liquidity providers respond by changing size at the best quotes rather than the spread \citep[same chap., sec.~2.2, p.~148]{alpha-laruelle}. In crashes liquidity suppliers withdraw: in 1987 buyers waited until they believed all sellers were satisfied, and some Nasdaq dealers took their phones off the hook \citep[chap.~``Bubbles, Crashes, and Circuit Breakers'', pp.~561 and~563]{alpha-harris}.

For large-tick instruments, Lehalle and Laruelle add a bounce term inside the same scaling by trading activity, $\psi_i=a+(b_1\sigma_i+b_2\,\eta\cdot\mathrm{tick})/\sqrt{N_i}$, as in \hyperref[meth:spread-volatility-regression]{the spread--volatility regression} (eq.~(2.1)). Here $\eta=N_c/(2N_a)$ is half the ratio of price continuations to alternations in traded prices: the more bid--ask bounce, the more alternations and the smaller $\eta$, which is one half when moves are random \citep[chap.~2, sec.~2.2, eq.~(2.1), pp.~144--145]{alpha-laruelle}. When the spread is stuck at one tick, liquidity providers who cannot step inside it increase the size at the best bid and offer, so a larger book plays the role of a narrower spread (p.~148). Interpretation: the tick is the constraint; bounce and book size are observed outcomes; matching rules (time priority or pro-rata allocation) govern who is filled within the larger queue (\hyperref[conv:cme-globex-trading]{CME Globex trading conventions}).

Futures often have economically large ticks, with bounce, reversals and larger displayed quantities than comparable equities; the mechanism matches these passages, which are not specific to futures. \emph{Futures evidence.} Dayri and Rosenbaum, on millisecond data for ten futures on four exchanges from 15~May to 31~December 2009, find all of them large-tick: the spread just before a trade was one tick for 72.7 (DAX) to 99.6 (S\&P~500) percent of trades.

The bounce statistic $\eta$, the same estimator $N_c/(2N_a)$, averaged 0.035 for the S\&P~500 future, 0.087 for the Euro Stoxx~50, 0.12 to 0.14 for the Schatz and Bund, and 0.23 to 0.28 for the DAX, Dow, EUR/USD, five-year Treasury and crude oil contracts, all below the one half of a random walk; they read a small $\eta$ as a large tick \citep[sec.~3.2, Table~1, pp.~18--19; $\eta$ estimator, p.~15]{dayri-rosenbaum-2013}. This supports ``often'' for liquid 2009 futures and the bounce. It contains no equities and no depth measure, so the comparison with equities and the larger displayed quantities remain unchecked.

\qffield{Pitfalls} Secondary factors are correlated with more than one primary factor and with each other, so an observed effect can have the opposite sign to the prediction (emerging industries attract volume), and tests need controls \citep[same chap., p.~314]{alpha-harris}. Volatility has a second, opposite effect through gambling and hedging interest, which raise activity \citep[same chap., pp.~316--317]{alpha-harris}. Trading activity is the best observable predictor of spreads, but Harris keeps it secondary to preserve the distinction between theoretical and observable determinants \citep[same chap., p.~312]{alpha-harris}. These are qualitative predictions; the importance ratings in Table~14-2 are Harris's judgement, not estimates. Harris's related factor is risk replication for hedgers. Among liquidity providers' responses to volatility, the source supports wider spreads and, in crashes, withdrawal, but not smaller displayed size or faster cancellations.

Lehalle and Laruelle report book size moving inversely to the spread over the day \citep[chap.~2, sec.~2.2, p.~148]{alpha-laruelle}; that is co-movement, not a measured response to volatility. Interpretation, consistent with Lehalle and Laruelle (pp.~145 and~148) but not stated by them: a low $\eta$ means more of the observed movement is microstructure noise, and a larger book replaces a narrower spread only for orders small enough to fit inside it.

\qffield{See also} \hyperref[res:adverse-selection-spread-component]{spread components}; \hyperref[res:glosten-milgrom-market-collapse]{market collapse} (a model of market failure); \hyperref[def:outside-spread]{outside spread} (the spread of value traders); \hyperref[meth:dealer-quote-adjustment]{dealer quote adjustment}; \hyperref[meth:spread-volatility-regression]{spread--volatility regression} (an empirical test of the volatility determinant); \hyperref[def:volatility-components]{fundamental and transitory volatility}.

\end{qfresult}

\begin{qfdefinition}[Outside spread: the prices at which value traders will trade]\label{def:outside-spread}
\qffield{Statement} A value trader's \emph{outside spread} is the interval between the prices at which he is willing to buy and to sell; value traders rarely quote these prices, so as not to reveal their value estimates or give free trading options \citep[chap.~``Value Traders'', sec.~16.2, p.~340]{alpha-harris}. Writing $\widehat V$ for his value estimate, he buys only below $b^o$ and sells only above $a^o$, and $s^o=a^o-b^o$. Interpretation: if every cost and risk of a one-unit trade could be summarised in an all-in half-cost $c>0$, the expected surplus would be $\widehat V-P-c$ from buying and $P-\widehat V-c$ from selling. He would trade only outside the no-trade band $[\widehat V-c,\widehat V+c]$, so $b^o=\widehat V-c$ and $a^o=\widehat V+c$. The spread widens with \citep[same chap., pp.~340--344]{alpha-harris}:
\begin{itemize}
\item adverse selection from news traders who know of value changes he has not yet learnt;
\item the winner's curse: the risk that his estimate is wrong, larger when values are hard to estimate and when many others estimate them, with losses from the market impact of trading on a wrong estimate;
\item the costs of research, which must be recovered also on the instruments analysed that offer no trade.
\end{itemize}

\qffield{Conditions} Value trading is profitable only when price differs from fundamental value, which happens in two ways. New information changes the value: news traders, by definition the first to learn it, profit and push the price to the new value. Or uninformed traders push the price away from value: value traders profit and push it back. Value traders therefore profit only by trading against uninformed traders' demands for liquidity \citep[same chap., sec.~16.1, pp.~338--339]{alpha-harris}. They do so directly, through limit orders or blocks, or indirectly through dealers whose inventories the uninformed demand has unbalanced \citep[same chap., secs.~16.1.1--16.1.2, pp.~339--340]{alpha-harris}.

Outside spreads are much wider than dealer spreads: value traders trade when order flow is one-sided and cannot learn their mistakes from it, take larger positions held until prices revert, so bear more inventory risk and financing cost, trade less volume over which to spread fixed costs, and earn at best half the outside spread when price returns to value, while a dealer often earns the whole spread on a round trip \citep[same chap., sec.~16.3, pp.~344--345]{alpha-harris}.

\qffield{Practice} Value traders buy when dealers lower their quotes to shed inventory taken from uninformed sellers, and so supply depth to the market and resiliency against uninformed price moves; dealers take larger positions where value traders follow the market closely \citep[same chap., pp.~340, 345--346]{alpha-harris}. When their value estimates differ enough, value traders' outside spreads overlap and they trade with each other \citep[same chap., p.~343]{alpha-harris}. Not buying does not oblige a trader to sell: uncertain estimates mean he should speculate only when he believes an instrument is very mispriced, and the larger the apparent mispricing, the more likely it is a mistake \citep[same chap., box, p.~343]{alpha-harris}. In continuous markets the winner's curse bites because traders trade at the prices they quote, and the market aggregates estimates sequentially; in a two-sided single-price auction it matters little \citep[same chap., pp.~342--343]{alpha-harris}.

\qffield{Pitfalls} The outside spread is not observed in quotes or trades (Interpretation: it is inferred from where value traders act). One classification lists leverage and borrowing among its financing costs; Harris names the cost of financing larger positions, without that detail \citep[same chap., p.~344]{alpha-harris}. A single half-cost $c$ understates the problem: adverse selection and the winner's curse are not per-unit costs but depend on what a trade at a given price reveals and on how many others are estimating the value, which is why Harris treats them as risks that widen the spread (Interpretation).

Potential sources of value-trader profit include uninformed liquidity demand and dealer inventory pressure. In Harris, information changes are where news traders profit; value traders gain there only by correcting news traders' under- or overreaction \citep[chap.~``Informed Traders and Market Efficiency'', Table~10-1, p.~230, and p.~239]{alpha-harris}. Dealer inventory pressure is the indirect channel of uninformed demand \citep[chap.~``Value Traders'', pp.~339--340]{alpha-harris}.

Dealer and outside spreads differ in valuation target, holding period, risks (inventory and client selection against valuation, financing and the winner's curse), scale and profit capture, as in the Conditions field. Dealer spreads stay tight even when fundamental investors require a wide outside spread, which is Harris's conclusion that outside spreads are much wider than dealer spreads (sec.~16.3, pp.~344--345). Client-selection risk for dealers is Harris's adverse selection from informed traders (\hyperref[def:dealer-quotes-order-flow]{dealers and their quotes}).

\qffield{See also} \hyperref[def:spread-measures]{spread measures} (contrast: displayed and transaction spreads); \hyperref[res:spread-determinants]{spread determinants}; \hyperref[meth:dealer-quote-adjustment]{dealer quote adjustment} (value traders restore dealers' inventories); \hyperref[def:fundamental-value]{fundamental value}; \hyperref[res:informed-trading-aggregation]{informed trading and aggregation}; \hyperref[def:winners-curse]{winner's curse}; \hyperref[def:informed-trader-types]{informed-trader types} (value against news traders).

\end{qfdefinition}
